\section{Data}
\label{sec:data}

\subsection{The \populace{} stack}
\label{sec:populace-stack}

\policyengineus{} runs its microsimulation over microdata built and
calibrated by \populace{}, PolicyEngine's open-source microdata stack: a
monorepo of five packages -- \texttt{populace-frame} (the weighted
sampling-frame kernel), \texttt{populace-fit} (weight-aware conditional
imputation), \texttt{populace-calibrate} (targets-to-weights calibration),
\texttt{populace-build} (country build plans and release gates), and
\texttt{populace-data} (the published population registry and release
publishing).
\evidence{/Users/maxghenis/PolicyEngine/populace/README.md, package table}
The predecessor US enhanced-CPS package, \texttt{policyengine-us-data}, is
archived; US microdata work is filed against \populace{} going forward.
\evidence{\texttt{gh repo list PolicyEngine}: policyengine-us-data shows
\texttt{public, archived}}

\subsection{Base construction}
\label{sec:base-construction}

The US base pool is assembled by a staged, checkpointed builder rather
than a single-process script: each build stage runs as its own subprocess
(so stage memory is released with the process), the per-target quantile
regression forest (QRF) imputation chain checkpoints between targets with
raw prior draws, exact target order, and both random-number streams
preserved, and the staged pipeline is proven bit-identical to the
monolithic builder at every stage boundary under a pinned single-thread
environment.
\evidence{\url{https://github.com/PolicyEngine/populace/pull/406} (merged
2026-07-13); equivalence proof at
/Users/maxghenis/PolicyEngine/populace/packages/populace-build/tests/test\_us\_puf\_support\_base\_equivalence.py}
Checkpoint reuse is guarded by a code-identity lock: the builder
fingerprints its own executable sources and dependency pins (a SHA-256
digest over the tool, all package sources, \texttt{pyproject.toml}, and
\texttt{uv.lock}) and records it as
\texttt{builder\_code\_identity.source\_sha256} in the run configuration,
so a resume against checkpoints written by different builder source fails
with a run-configuration mismatch rather than silently mixing code
versions.
\evidence{/Users/maxghenis/PolicyEngine/populace/tools/build\_us\_puf\_support\_base.py,
\texttt{\_builder\_code\_identity()}; the refusal is pinned by
\texttt{test\_outer\_stage\_resume\_rejects\_changed\_builder\_code} in
packages/populace-build/tests/test\_us\_puf\_support\_base\_builder.py}
QRF fitting is memory-bounded for near-discrete targets: when a target
variable has 2--32 distinct values, \texttt{max\_samples\_leaf} is bounded
at 4{,}096 rather than storing full leaf populations, which exposes each
leaf quantile to multinomial sampling noise while preventing the dense
leaf-storage growth such targets otherwise cause.
\evidence{/Users/maxghenis/PolicyEngine/populace/packages/populace-fit/src/populace/fit/qrf.py,
\texttt{\_DISCRETE\_Y\_UNIQUE\_MAX} (32) and \texttt{\_DISCRETE\_Y\_LEAF\_BOUND}
(4096); \url{https://github.com/PolicyEngine/populace/pull/415}}

The base pool represents each source tax unit on two support channels: an
\asec{} channel holding survey-measured values and a
\texttt{puf\_tax\_detail} channel holding IRS Public Use File (PUF) donor
detail. Tax-detail columns that only the PUF measures are populated
exclusively on the PUF channel -- exactly zero on the \asec{} channel by
design -- while survey-measured columns are carried on both channels (a
unit's PUF clone rows carry its survey predictors). \populace{}
deliberately does not reproduce the retired pipeline's both-half override,
which overwrote measured survey values with PUF imputations.
\evidence{/Users/maxghenis/PolicyEngine/populace/packages/populace-build/src/populace/build/us/source\_stages.json,
PUF-detail stage notes (``places this PUF-only detail on its dedicated PUF
support channel \ldots rather than reproducing the retired both-half
override''; ``deliberately preserves measured ASEC operations income and
replaces only the PUF support channel''); channel constants in
packages/populace-build/src/populace/build/us\_runtime/capital\_gain\_details.py;
clone-predictor behavior in
\url{https://github.com/PolicyEngine/populace/pull/425}}

Where the frozen census extracts \populace{} builds from lack a column the
retired pipeline measured, the column is restored from pinned official
Census \asec{} public-use archives rather than re-imputed. Each income
year's person archive is pinned by zip and member SHA-256, byte size, and
CRC plus row and audit pins; restoration is an exact \texttt{PERIDNUM}
join within each income year and never predicts a source value. The
mapping from a \texttt{census\_cps} income year to the \asec{} survey year
one year later is verified empirically per year: each income-year cohort
has 100.0\% \texttt{PERIDNUM} coverage in its survey-year-plus-one archive
and only roughly 33\% (one rotation group) in adjacent vintages, so a
wrong archive pin fails the exact join loudly instead of silently
degrading. Two columns are restored this way: weeks unemployed (\asec{}
\texttt{LKWEEKS}) and educational assistance (\asec{} \texttt{ED\_VAL}).
\evidence{\url{https://github.com/PolicyEngine/populace/pull/421} (merged
2026-07-13);
/Users/maxghenis/PolicyEngine/populace/packages/populace-build/src/populace/build/us\_runtime/education\_assistance\_source.py,
module docstring and \texttt{AsecEducationArchive} pins}

\subsection{The input-coverage contract}
\label{sec:input-coverage}

A microsimulation input that is silently absent produces a specific
failure mode: a reform that changes policy on that input scores exactly
zero, with no error. \populace{}'s US releases encode the guarantee
against this failure directly, and the gate that enforces it was
introduced after exactly such a failure -- an SSI resource-limit reform
(\$10k/\$20k asset limits) scored \$0 on a certified release because the
three countable-resource asset inputs were absent from the export. The
coverage gate shipped deliberately red on the then-current artifacts, with
the asset columns as unconditional requirements.
\evidence{\url{https://github.com/PolicyEngine/populace/pull/369};
\url{https://github.com/PolicyEngine/populace/issues/368}}

The contract is a versioned, in-repo register covering every input column
the pinned reference eCPS export populates: each column is either
\texttt{required} -- it must be present in the release with non-default
signal -- or a \texttt{reviewed\_exclusion} carrying a written reason and
a tracking issue. The release gate fails a required column that is absent
\emph{or} present-but-degenerate (every value at the engine default), and,
with cannot-rot semantics, fails a reviewed exclusion whose column has
since been restored; a companion reform-coverage smoke simulates a bound
reform against the written artifact at release time, so reform
responsiveness is demonstrated rather than inferred. Register edits are
themselves guarded: a fail-closed sync guard rejects exclusions for
columns the build now populates (during the Build~M campaign it correctly
rejected an exemption proposed for a restored input; the measurement, not
the register, was fixed).
\evidence{/Users/maxghenis/PolicyEngine/populace/packages/populace-build/src/populace/build/us/release\_input\_coverage\_manifest.json,
\texttt{description} field; regenerated by
tools/build\_us\_release\_input\_coverage\_manifest.py; the
rejected-exemption episode in
\url{https://github.com/PolicyEngine/populace/pull/428}}

The coverage campaign then moved the contract from mostly-excluded to
mostly-required: at Build~J certification (2026-07-10) the US contract
stood at 70 required columns against 88 reviewed exclusions; the contract
on \texttt{main} at this revision requires 158 of 166 columns, with 8
reviewed exclusions -- the contract Build~M builds under. The resulting
guarantee is the one that matters for reform scoring: a reform touching a
required-and-populated input cannot silently score zero because the input
is missing, and the columns for which that guarantee is \emph{not} made
are enumerated with reasons rather than absent without notice.
\evidence{Build~J counts:
\url{https://github.com/PolicyEngine/populace/issues/368}, certification
verdict comment (``65$\to$70 required / 93$\to$88 exclusions''); current
counts:
/Users/maxghenis/PolicyEngine/populace/packages/populace-build/src/populace/build/us/release\_input\_coverage\_manifest.json,
\texttt{counts}: 158 required, 8 reviewed\_exclusion, 166 total}

\subsection{Calibration to administrative totals}
\label{sec:calibration}

The microdata is calibrated to a target surface of administrative and
official aggregates: IRS \soi{}, Census, SSA, CMS, SNAP, and TANF program
totals, and \cbo{}/\jct{} figures, resolved from an external PolicyEngine
Ledger consumer-facts feed rather than hardcoded into the calibration
config.
\evidence{/Users/maxghenis/PolicyEngine/populace/packages/populace-build/src/populace/build/us/fiscal\_target\_references.json,
\texttt{description} field: "Observed values must be resolved from external
PolicyEngine Ledger consumer facts; generated IRS, Census, SSA, CMS, SNAP,
TANF, and CBO targets are compiled directly from the Ledger fact stream."}
Per-target calibration fit is recorded in \texttt{calibration\_diagnostics.json}
at release build time: each target row carries the official value
(\texttt{target}), \policyengineus{}'s calibrated estimate
(\texttt{final\_estimate}), a \texttt{relative\_error}, and a
\texttt{within\_tolerance} flag, gated at release time against a named
tolerance per target family -- for example, federal income tax amount at
5\% and \soi{}/SSA-keyed EITC, CTC, ACTC, PTC, and taxable Social Security
targets at 15\%, and SALT at 10\%.
\evidence{/Users/maxghenis/PolicyEngine/populace/packages/populace-data/src/populace/data/contract.py,
\texttt{\_US\_CRITICAL\_TARGET\_FIT\_REQUIREMENTS}; diagnostics schema at
/Users/maxghenis/PolicyEngine/populace/packages/populace-calibrate/src/populace/calibrate/diagnostics.py}

For the currently certified US release -- Build~J, release id
\texttt{populace-us-2024-buildj-sparse-rmloss100-75d5add-20260710T094201Z},
certified 2026-07-10 against \policyengineus{} 1.764.6 -- the
certification verdict records a final calibration loss of 0.030843 and
89.0\% of the registry's 5{,}510 compiled targets within 10\% relative
error, with every release gate enforced and no bypass flags
(Section~\ref{sec:release-gates}).
\evidence{\url{https://github.com/PolicyEngine/populace/issues/368},
certification verdict comment (2026-07-10): final loss 0.030843;
within-10\% 0.8900; registry \texttt{091ace02f962}, 5{,}510 compiled
specs; \texttt{release\_gates.passed=True}, \texttt{failures=[]}}

\todo{For v1 submission, the per-target table behind these headline
figures (actual relative\_error values by target family) must be fetched
from the release bundle's \texttt{calibration\_diagnostics.json} on the
Hugging Face dataset \texttt{policyengine/populace-us} for the release
pinned in Section~\ref{sec:disclosures}; the headline figures above are
citable to the certification verdict, but per-target figures must come
from the published artifact, not from memory or a prior release.}

\todo{Build~M scorecard slot. Build~M is in certification at the time of
this revision (2026-07-15) and its numbers must not be cited before its
certification verdict publishes. When it does, record here: the release
id, final calibration loss, within-10\% share, export-mass parity, and
reform-probe results, each copied from the verdict. State two
comparability caveats alongside: (1) Build~M's registry compiles more
targets than Build~J's 5{,}510 (the take-up campaign added SSA
recipient-count target families), so within-10\% shares are not directly
comparable across builds -- more targets is a harder calibration problem;
and (2) Build~M builds under the 158/8 input-coverage contract
(Section~\ref{sec:input-coverage}) versus Build~J's 70/88, so its input
surface is materially wider. Until that verdict, Build~J remains the
certified default and the only citable scorecard.}

\subsection{Take-up}
\label{sec:takeup}

Benefit take-up enters the data as seeded stochastic flags (for example
\texttt{takes\_up\_ssi\_if\_eligible}), restoring parity with the retired
pipeline's stochastic take-up inputs across programs.
\evidence{\url{https://github.com/PolicyEngine/populace/issues/312}}
For SSI the flags are count-calibrated against official recipient counts
rather than drawn at a fixed rate, with a saturation fallback: the
assigner draws a population-wide propensity equal to the ratio of the
official target to modeled candidate capacity (capped at one) and
count-calibrates candidates, and when every age band saturates (candidate
capacity below the official target), the reform-domain propensity falls
back to the observed take-up rate among modeled candidates rather than
degenerating to certainty for the entire pool; current-law recipiency is
untouched in both regimes.
\evidence{\url{https://github.com/PolicyEngine/populace/pull/423} (merged
2026-07-14)}

Because calibration (Section~\ref{sec:calibration}) re-weights the pool
after flags are assigned, flags assigned under pre-solve weights need not
remain consistent with the returned weights. The Build~M certification
campaign proved two structural defects in the original reconciliation
loop: the registry's national SSI-recipients target and the take-up
stage's SSA age-band sum are different SSA measures roughly 115{,}000
apart, so enforcing both simultaneously on the modeled candidate mass gave
the fixed-point iteration no fixed point; and each pass gated frozen flags
against weights that had since drifted (a stale pair). The repaired loop
aligns the band goals to SSA age \emph{shares} applied to the registry
national total (extracted fail-closed from the compiled specs), and each
pass re-assigns flags under the current weights, replays the dependent
health tail, re-materializes the affected targets, and refits the
calibration; the loop exits only on a fresh pair -- one further assignment
plus health-tail replay under the returned weights, with the gates run on
that self-consistent pair and per-band deltas recorded.
\evidence{\url{https://github.com/PolicyEngine/populace/pull/430} (merged
2026-07-14)}

The fresh-pair exit's swap delta -- the change in selected recipient mass
when flags are re-assigned under the returned weights -- measures the
calibration solve's equilibrium residual on the SSI-recipient target
family, not assignment granularity: the solve balances that family against
the rest of the target surface (5{,}692 other targets in the diagnosed
run) and can land short of the national total (about 5.7\% in that run),
and the fresh assignment then restores recipiency to the official counts.
The delta is therefore recorded in the reconciliation manifest with the
assignment-granularity sum as a reference quantity, and gated only against
runaway, at one tenth of the fresh national total -- a solve that
effectively abandons the SSI-recipient family still fails closed, with the
delta in the terminal error.
\evidence{\url{https://github.com/PolicyEngine/populace/pull/431} (merged
2026-07-15)}

The take-up architecture carries a documented limitation: the eligibility
model currently produces roughly 4 million persons of modeled SSI
candidate mass against the registry's national recipients target of
7{,}404{,}820 (SSA, December 2024), with candidate capacity below the
official target in every SSA age band. Calibration cannot manufacture
eligible households that the eligibility model rejects, so recipient-count
calibration concentrates weight on the candidates that do exist; the
undercount is tracked as a modeling-side defect, and any v1 claim about
SSI reform aggregates must state it.
\evidence{\url{https://github.com/PolicyEngine/populace/issues/424}
(the 7{,}404{,}820 SSA December-2024 total and the per-band shortfall);
\url{https://github.com/PolicyEngine/populace/pull/430}
(``\textasciitilde{}4M of modeled candidate mass'')}

\subsection{Release provenance}
\label{sec:release-provenance}

Each \populace{} US release publishes as an immutable, versioned bundle to
the Hugging Face dataset \texttt{policyengine/populace-us}
(release ids of the form \texttt{populace-us-<year>-<gitsha>-<date>}), gated
on a release contract requiring a build manifest (git commit, non-dirty
working tree, dataset and calibration artifact SHA-256 hashes), a release
manifest (compatible core/model package version ranges, artifact locations),
calibration diagnostics, and (for US) a source-coverage manifest. Release
metadata additionally includes a TRACE Transparent Research Object (a
JSON-LD provenance record pinning build inputs and outputs by hash) for at
least one build.
\evidence{/Users/maxghenis/PolicyEngine/populace/packages/populace-data/src/populace/data/contract.py
(REQUIRED\_RELEASE\_FILES, schema versions); example TRO at
/Users/maxghenis/PolicyEngine/populace/packages/populace-data/build/us/populace\_us\_2024.trace.tro.jsonld,
build id \texttt{populace-us-2024-9f1260b-20260611}} Publishing is refused
if the release's out-of-sample reform validation (Section~\ref{sec:validation})
has not actually been simulated (the \texttt{out\_of\_sample\_simulated}
flag).
\evidence{/Users/maxghenis/PolicyEngine/populace/packages/populace-data/src/populace/data/release.py}

The currently certified default is the Build~J bundle: the dataset's
\texttt{latest.json} points to
\texttt{populace-us-2024-buildj-sparse-rmloss100-75d5add-20260710T094201Z},
published 2026-07-10 as a tracked, tagged publish and certified for
consumption through \policyenginepy{}.
\evidence{\url{https://github.com/PolicyEngine/populace/issues/368},
``SHIPPED as the certified US default'' comment: ``latest.json ->
populace-us-2024-buildj-sparse-rmloss100-75d5add-20260710T094201Z (tracked
publish, tagged)''; certification wiring in
PolicyEngine/policyengine.py\#470}

\todo{At submission time, re-verify this pin directly against the Hugging
Face dataset's published \texttt{latest.json} -- Build~M is in
certification as of this revision and the default may flip -- and record
the final pin in Section~\ref{sec:disclosures}. The informal identifier
``4.18.7 (dense)'' from the seed notes is superseded by the release ids
above and must not be cited.}
